% Chapter 18: The Cap Set Problem
\let\LocalMacro\relax

\chapter{The Cap Set Problem}

\section{The Problem}

\subsection{Informal}
Imagine a big grid of points in space where the arithmetic wraps around like a clock. You want to pick as many points as possible, but with one strict rule: no three of your chosen points can be equally spaced in a line. How many points can you pick before you are forced to have three in a row? This problem stumped mathematicians for decades, despite being easy to state.

\subsection{Formal}
\begin{conjecture}[Frankl-Graham-Rödl, 1987]
$r_3(n) = o(3^n)$.
\end{conjecture}

The trivial bound is $3^n$. Meshulam (1995) proved $r_3(n) \le 2 \cdot 3^n / n$, using Fourier analysis. The best pre-2016 bound was $r_3(n) \le 3^n / n^{1+c}$ for a small constant $c > 0$.

\begin{historicalbox}
The cap set problem is the $\mathbb{F}_3$-analogue of Roth's theorem on arithmetic progressions in the integers (1953). While Roth's theorem was proved using Fourier analysis on $\mathbb{Z}$, the finite-field version resisted analogous techniques for decades.
\end{historicalbox}

\subsection{Why It Matters}
The cap set problem is a testing ground for new mathematical techniques. The slice rank method used to solve it has since been applied to many other problems in combinatorics and theoretical computer science. Understanding how large structured sets can be has applications to coding theory and the study of error-correcting codes.

\subsection{What Was Known}
Roth proved in 1953 that subsets of the integers without three-term arithmetic progressions must be thin. Meshulam extended this to three-dimensional grid arithmetic in 1995, proving that cap sets must be significantly smaller than the full grid. But the best known bound was still far from the conjectured answer. The polynomial method, newly adapted to combinatorics by Croot, Lev, and Pach in 2016, provided the breakthrough.

\section{Mathematical Theory}

\subsection{Prerequisites}
This chapter assumes familiarity with:
\begin{itemize}
\item Basic combinatorics (binomial coefficients, extremal combinatorics)
\item Linear algebra (rank of matrices and tensors)
\item Finite fields ($\mathbb{F}_3$, arithmetic in characteristic 3)
\end{itemize}

\subsection{The Polynomial Method}

The polynomial method is a versatile combinatorial tool with a simple core idea: to show that a set $A$ is small, find a nonzero polynomial that vanishes on all of $A$. If the polynomial has low degree, it cannot have too many zeros, so $A$ must be small.

\begin{definition}[Vanishing Polynomial]
A polynomial $P \in \mathbb{F}_q[x_1, \ldots, x_n]$ \emph{vanishes} on a set $A \subset \mathbb{F}_q^n$ if $P(a) = 0$ for all $a \in A$.
\end{definition}

The key fact connecting polynomials to finite fields is:

\begin{fact}[Schwartz-Zippel]
A nonzero polynomial of degree $d$ over $\mathbb{F}_q^n$ has at most $d \cdot q^{n-1}$ zeros.
\end{fact}

Dvir's 2008 proof of the Kakeya conjecture over finite fields demonstrated the power of this approach: by constructing a low-degree polynomial vanishing on a Kakeya set, he showed such sets must be large. The cap set problem asks whether a similar strategy can beat Meshulam's Fourier-analytic bound.

\subsection{Tensor Rank and Slice Rank}

The breakthrough of Croot, Lev, Pach, and Ellenberg--Gijswijt reformulates the cap set problem in terms of \emph{tensor rank}, a generalization of matrix rank.

\begin{definition}[Tensor]
A 3-tensor $T \in V_1 \otimes V_2 \otimes V_3$ over a field $\mathbb{F}$ is a multilinear map $T: V_1^* \times V_2^* \times V_3^* \to \mathbb{F}$.
\end{definition}

\begin{definition}[Slice Rank]
The \emph{slice rank} of a 3-tensor $T \in V_1 \otimes V_2 \otimes V_3$ is the minimum $r$ such that $T$ can be written as a sum of $r$ terms, each of which is a product of a vector in one factor with a bilinear form on the other two:
\[
T = \sum_{i=1}^{r} v_i \otimes B_i
\]
where each $v_i$ lives in one of $V_1, V_2, V_3$ and $B_i$ lives in the tensor product of the other two.
\end{definition}

Slice rank is always at most the tensor rank, and for diagonal tensors (those supported only on the ``diagonal'' $x = y = z$), slice rank equals the size of the diagonal support.

\subsection{The Cap Set Tensor}

For the cap set problem over $\mathbb{F}_3^n$, define the tensor:
\[
T(x,y,z) = \begin{cases} 1 & \text{if } x + y + z = 0 \\ 0 & \text{otherwise} \end{cases}
\]
A subset $A \subset \mathbb{F}_3^n$ is a cap set if and only if $T|_{A \times A \times A}$ is diagonal (i.e., $T(x,y,z) \neq 0$ only when $x = y = z$). The slice rank of $T$ then upper-bounds $|A|$:
\[
|A| \le \text{slice rank}(T).
\]

\begin{highlightbox}
The key insight: the cap set condition is exactly the condition that the restriction of $T$ to $A$ is diagonal. Since diagonal tensors have slice rank equal to the size of the diagonal, any cap set has size at most the slice rank of $T$.
\end{highlightbox}

\subsection{Counting Monomials}

The slice rank of $T$ is bounded by counting low-degree monomials. Consider the vector space $\mathcal{P}_{<2n/3}$ of polynomials in $\mathbb{F}_3[x_1, \ldots, x_n]$ of degree less than $2n/3$ (with each variable exponent in $\{0, 1, 2\}$). Using Stirling's approximation:

\begin{theorem}[Monomial Count]
$|\mathcal{P}_{<2n/3}| \le 2.756^n$.
\end{theorem}

The slice rank of $T$ is at most $3 \cdot |\mathcal{P}_{<2n/3}|$ (the factor 3 accounts for the three tensor factors), giving:
\[
r_3(n) \le 3 \cdot 2.756^n.
\]

\section{The Solution}

\subsection{Ellenberg-Gijswijt (2016)}

\begin{theorem}[Ellenberg-Gijswijt \cite{ellenberggijswijt2017}]
$r_3(n) \le 2.756^n$ for sufficiently large $n$.
\end{theorem}

The key idea is the \emph{slice rank} method, a variant of the polynomial method:

\begin{enumerate}
\item \textbf{Tensor formulation}: Represent the cap set condition by a tensor $T: \mathbb{F}_3^n \times \mathbb{F}_3^n \times \mathbb{F}_3^n \to \{0,1\}$ defined by $T(x,y,z) = 1$ if $x + y + z = 0$ and $0$ otherwise. A cap set is a set $A$ such that $T|_{A \times A \times A}$ is diagonal (no off-diagonal entries).

\item \textbf{Slice rank bound}: The slice rank of $T$ is at most $3 \cdot M$, where $M$ is the number of monomials $x_1^{e_1} \cdots x_n^{e_n}$ with $e_i \in \{0,1,2\}$ and $\sum e_i < 2n/3$.

\item \textbf{Stirling's formula}: $M \le 2.756^n$, giving the bound $|A| \le \text{slice rank}(T) \le 3 \cdot 2.756^n$.
\end{enumerate}

\begin{highlightbox}
The slice rank method is a powerful generalization of the polynomial method: instead of vanishing of a polynomial, one uses the rank of a tensor. It has since been applied to many problems including the sunflower lemma and the Erd\H{o}s-Ginzburg-Ziv theorem.
\end{highlightbox}

\section{Impact}
\begin{itemize}
\item \textbf{Polynomial method 2.0}: Slice rank opened a new chapter in the polynomial method, extending it from geometry (Dvir's Kakeya) to combinatorics.
\item \textbf{Roth's theorem over $\mathbb{F}_3$}: The bound $2.756^n$ is exponentially better than previous methods.
\item \textbf{Connections to TCS}: The slice rank is related to the non-deterministic communication complexity of tensor operations.
\end{itemize}

\section{Exercises}

\begin{exercise}[Basic]
Show that a cap set in $\mathbb{F}_3^n$ has no three points on a line. Compute the maximum cap set size for $n=2$.
\end{exercise}

\begin{exercise}[Intermediate]
Explain how the slice rank of a tensor generalizes the rank of a matrix. Why does slice rank bound cap sets?
\end{exercise}

\section{Timeline}

\begin{tabularx}{\linewidth}{lX}
\toprule
\textbf{Year} & \textbf{Event} \\ \midrule
1953 & Roth proves no 3-AP in dense subsets of $\mathbb{Z}$ \\
1995 & Meshulam proves $r_3(n) \le 2 \cdot 3^n / n$ via Fourier analysis \\
2016 & Croot-Lev-Pach use polynomial method for $\mathbb{Z}_4^n$ \\
2016 & Ellenberg-Gijswijt solve cap set with slice rank \\
2022 & Huh wins Fields Medal (log-concavity in combinatorics) \\
\bottomrule
\end{tabularx}

\section{Citations}

\begin{enumerate}
\item Ellenberg, J., \& Gijswijt, D. (2017). ``On large subsets of $\mathbb{F}_3^n$ with no three-term arithmetic progression.'' Annals of Mathematics, 185(1), 339--343.
\item Tao, T. (2016). ``A symmetric formulation of the Croot-Lev-Pach-Ellenberg-Gijswijt capset bound.'' Blog post.
\item Croot, E., Lev, V., \& Pach, P. P. (2017). ``Progression-free sets in $\mathbb{Z}_4^n$ are exponentially small.'' Annals of Mathematics, 185(1), 331--337.
\item Meshulam, R. (1995). ``On subsets of finite abelian groups with no 3-term arithmetic progressions.'' Journal of Combinatorial Theory, Series A, 71(1), 168--174.
\end{enumerate}
